


\subsection{Fundamentals of Credit risk CR-1}

\subsubsection{基本概念}
\begin{itemize}
\item 信用风险是由于对方inability、unwillingness，或nontimeliness履行财务义务而导致经济损失的可能性。
\item Borrower、obligor、counterparty和bond issuer是接收信贷的一方。
\item Lender、creditor和obligee是提供信贷的一方。
\item Inability、nontimeliness to Pay
    \begin{itemize}
        \item 企业的现金流不足以偿还为快速扩张的贷款。 
        \item 产品或服务过时导致收入不足，无法支付运营和融资成本。
        \item 意外事件(如自然灾害、疫情)导致企业陷入财务困境。
        \item 如果款项未能及时偿还，债权人可能会错失利息收入或产生营运资本融资费用。
    \end{itemize}
\item Unwillingness to Pay
    \begin{itemize}
    \item 合同争议导致的商业纠纷 (Commercial disputes) 是最常见的情况。
    \item 如果纠纷进入法律诉讼程序，且贷方在诉讼中获胜 (the lender prevails)，那么债务就可以得到偿还。
    \item 不典型的情况——主权国家(Sovereign state)不履行其国际债务义务，或强行将其外币债务转换为本国货币。
    \end{itemize}
\end{itemize}

\subsubsection{Insolvency、Default和Bankruptcy的区别}
\begin{itemize}
\item Insolvency（资不抵债）:【Negative equity：Liabilities＞资产 Assets】

Insolvency ≠ Bankruptcy【把"资不抵债"看作是"破产"的同义词是错误的】

 Insolvency ≠  default，只需在下一次支付资金前筹集到资金。

\item Default（违约）: 【对手方未能履行合同义务：如未付款】

 Default通常是由于资不抵债或流动性不足（Illiquidity）。

\item Bankruptcy（破产）: 违约实体可以通过申请破产寻求法律保护。

\begin{itemize}

\item 法院会审查违约实体的财务状况，与管理层、债权人以及股东进行谈判，通过出售资产和/或与贷款人重新谈判融资安排，试图让违约实体继续经营。

\item 如果业务无法重组，则企业解散。
\end{itemize}

【破产法第7章主要涉及解散dissolution/清算liquidation；第11章主要涉及重组restructuring/reorganization】

\end{itemize}





\subsubsection{产生信用风险的交易类型}

\begin{center}
    \renewcommand{\arraystretch}{1.3}  % 增加行高
    \adjustbox{max width=\textwidth}{
    \begin{tabular}{|c|c|c|}
    \hline
    \makecell
    {\textbf{信贷类型}\\\textbf{Credit Type}} &
    \makecell{\textbf{损失来源}\\\textbf{Losses Result From}} &
    \makecell{\textbf{损失类别}\\\textbf{Loss Type}} \\
    \hline

         \makecell[c]{Loaned money\\贷款} 
        & \makecell[c]
        {• Nonrepayment 不还\\
        • Slow repayment 慢还\\
        • Dispute 争议\\
        •Enforcement 强制执行} 
        & \makecell[c]
        {• Face amount 本金\\
        • Interest 利息\\
        • Time value of money 货币时间价值\\
        • Frictional costs 摩擦成本} \\
    \cline{1-3}
         \makecell[c]
        {Lease obligation\\租赁义务} 
        & \makecell[c]{• Nonpayment 未付款} 
        & \makecell[c]
        {• Recovery of asset 回收租赁资产\\
        • Remarketing costs 再营销成本\\
        • Difference in conditions 条件变化损失} \\
    \cline{1-3}
         \makecell[c]{Receivables\\应收账款} 
        & \makecell[c]{• Nonpayment of goods delivered \\ or service performed \\客户未支付\\已交付的商品或服务} 
        & \makecell[c]{• Face amount}\\
    \cline{1-3}
        \makecell[c]{Prepayment for goods or services\\商品或服务的预付款} 
        & \makecell[c]
        {• Nondelivery 供方未交付\\
        • Performance not as contracted\\ 
        交付不符合合同规定\\
        • Slow delivery 交付延迟\\
        • Dispute/enforcement} 
        & \makecell[c]
        {• Replacement cost\\
        购买替代品的成本\\
        • Incremental cost\\
        额外的运营成本\\
        • Time value of money\\
        • Frictional costs} \\
    \cline{1-3}
        \makecell[c]{Deposits\\存款} 
        &\makecell[c]{• Nonrepayment}
        &\makecell[c]{• Face amount\\• Time value of money} \\
    \cline{1-3}
        \makecell[c]{Claim or contingent claim on asset\\对资产的索赔/或有索赔} 
        & \makecell[c]{• Nonrepayment/Noncollection\\• Slow repayment collection\\• Dispute/enforcement} 
        & \makecell[c]{• Face amount\\• Time value of money\\• Frictional costs} \\
    \hline
        \makecell[c]{Derivative\\衍生品} 
        & \makecell[c]{• Default of third party\\第三方违约} 
        & \makecell[c]{• Replacement cost\\（mark-to-market value）} \\
    \hline
    \end{tabular}
    }
    \end{center}

\subsubsection{存在信用风险的实体}
\begin{itemize}
\item A. Financial Institutions 金融机构
    \begin{itemize}
    \item Banks 银行
        \begin{itemize}
        \item 因发放贷款拥有庞大的信贷组合。
        \item 信贷风险来自贷款和信贷额度（信贷额度或信用额度）、基于资产的借贷（如回购）以及衍生工具。
        \item 通过复杂的风险管理实践和资本市场对冲敞口管理信用风险。
        \end{itemize}
    \item Asset Managers 资产管理公司
        \begin{itemize}
        \item 例如贝莱德 (Blackrock) 或先锋集团 (Vanguard Group) 
        \item 代表个人和机构投资大量资金。
        \item 信用风险来源于投资，尤其是债务工具。
        \end{itemize}
    \item  Hedge Fund对冲基金
        \begin{itemize}
        \item 投资者有更大的风险偏好，但要求高回报以补偿这种风险。
        \item 往往参与风险较高的投资和策略，如不良债务和卖空，从而从信用恶化中获利（generate profits from credit deterioration）。
        \item 一方面承担比较高的信用风险， 另一方面通过高信用风险获利。
        \end{itemize}
    \item Insurance Companies 保险公司
        \begin{itemize}
        \item 保险公司在承保活动、投资组合和再保险可收款方面面临信用风险。
        \end{itemize}
    \item Pension Funds 养老金基金
        \begin{itemize}
        \item 基金的一大部分投资于信用风险资产。
        \end{itemize}
    \end{itemize}

\item B. Corporations【主要指非金融公司】
    \begin{itemize}
    \item 信用风险来源
        \begin{itemize}
        \item 应收账款（accounts receivables）【应收账款相当于给客户提供短期贷款】
        \item 短期投资和银行存款
        \item 衍生品
        \item 供应商融资（Vendor financing）：如 IBM 全球融资等公司提供供应商融资，承担与银行类似的信用风险。
        \item 供应链风险（supply chain）：关键供应商的破产可能会带来风险。
        \end{itemize}
    \end{itemize}

\item C. Individuals 个人
    \begin{itemize}
    \item 家庭信贷风险
    \item 投资活动
    \item 存款
    \item 个人预付款项（如预付给承包商的款项）等
    \end{itemize}
\end{itemize}


\subsubsection{管理信贷风险的动机}




\begin{itemize}
\item Survival 生存

金融机构尤其需要关注信用风险管理，大额损失可能会导致破产；非金融企业的信用损失也可能引发破产。

\item Profitability 盈利能力

减少损失即增加利润——这对于低利润率的企业非常关键。

\item Return on equity 股本回报率

企业需要足够的股本资本和谨慎的风险管理，在债务和股权之间找到合适的平衡点，以实现股权回报的最大化。
\end{itemize}


\subsection{Governance CR-2}
实践中，尽管个别交易可能会使得公司面临大量损失，但通常不会导致公司破产。公司破产往往是积累了大量的有毒交易。这通常源于风险管理框架薄弱(poor risk management framework)和公司治理失效(corporate governance failure)。

本考点主要强调在大型公司中，良好的信用风险管理框架和公司治理的重要性。


\subsubsection{Three Lines of Defense三道防线}

\begin{itemize}

    \item  第一道防线：负责并管理风险的\textbf{\textcolor{red} {业务负责人}}（own and manage risk）
    \item  第二道防线：监督风险，建立政策和程序， 对第一道防线进行管理监督（oversight for the first line）的 \textbf{\textcolor{red} {风险管理、合规和法律部门}}（enter prise risk
    management， compliance， and legal）
    
    \item  第三道防线：独立（provides independent assurance）管理和监控第一和第二道防线的风险的\textbf{\textcolor{red} {内部审计、外部审计和特别审计委员会}}（Internal and external auditors and audit committees）
\end{itemize}

【在操作风险和流动性风险中也有三道防线。大原则是一致的，主要是第一道防线（业务负责人）不同。例如，在信用风险中第一道防线是信贷客户经理。流动性风险中第一道防线是司库部门】
    
 

\subsubsection{最佳治理体系的四个关键原则：}

\begin{itemize}
\item A. Guidelines指导方针： 【明确的指导方针，管理关于信用风险交易的审批。】
    \begin{itemize}
    \item 指导方针是解释交易达成前必须遵守的规则的一系列文件，也叫信贷政策（credit policies） 或风险管理标准（risk management standards）等。
    \item 特征：
        \begin{itemize}
        \item Understandable易于理解：语言清晰简单，易于理解，避免使用过于法律化的语言
        \item Concise简明扼要
        \item Precise准确
        \item Accessible易于获取
        \end{itemize}
    \item Creation and Approval Process创建和审批流程
        \begin{itemize}
        \item 必须由高级管理人员（例如CRO、CFO）发起， 并由輩事会批准。
        \end{itemize}
    \item Promulgating and Maintaining Guidelines颁布和维护
        \begin{itemize}
        \item 首席风险官办公室负责指导方针的起草、批准、推广和维护。
        \item  指导方针需要定期更新和审查， 与业务发展同步。
        \end{itemize}
    \item Content of Guidelines 指导方针内容
        \begin{itemize}
        \item 指导方针的目的
        \item 定义交易关键参数的方法
        \item 交易审批流程和授权委托
        \item 新产品和市场的处理流程
        \item 审查和更新的流程
        \item 不遵守指导方针的后果
        \end{itemize}
    \item Breach（违反） of Guidelines
        \begin{itemize}
        \item 违反Guidelines是一种严重行为， 可能导致立即解雇。
        \item 也可以考虑一些特殊情况（如因外汇波动违反信用额度），设置一些特例。
        \end{itemize}
    \end{itemize}

\item B. Skills: 【将权限下放至具备合适技能的委员会和个人】
    \begin{itemize}
    \item 理想情况下， 所有交易都应由公司的高级管理人员批准，但实际工作中不太现实。这就涉及到授权委派问题。
    \item 风险管理与业务单位的冲突：
        \begin{itemize}
        \item 业务单位的优先事项： 赚钱（第一）， 不亏钱（第二），希望最大自由度。
        \item 风险管理的优先事项： 不亏钱（第一）， 赚钱（第二），需要限制自由度。
        \end{itemize}
    \item Delegation of Authority权力下放的步骤
        \begin{itemize}
        \item Step 1: 从风险管理的角度，为每笔交易分配基本参数
        \item  Step 2: 基于这些参数委派批准权限。
        \end{itemize}
    \item Defining risk parameters 定义风险参数
        \begin{itemize}
        \item Amount of exposure敞口：估计公司可能损失的最大金额。
        \item  Credit quality 信用质量：必须制定一个标准总结每个交易对手的信用状况。
        \item  Tenor (length) of the exposure 敞口期限：面临潜在损失的期间
        \end{itemize}
    \item Delegation of Authority 授权委托的规则：
        \begin{itemize}
        \item 交易风险越高，所需的审批级别越高。
        \item 高风险、低信用质量或长期交易，需要高级管理层的关注，
        必须由具备资格的委员会审批。
        \item 小额、短期且信用质量高的交易可以在较低级别批准。
        \item 简单明了的交易可以由个人批准。
        \end{itemize}
    \begin{figure}[H]    % 强制图片在当前位置
        \centering
        \includegraphics[width=0.8\textwidth]{process.png}
        \caption{process}
        \label{fig:process}
    \end{figure}
    \end{itemize}

\item C. Limits：【设定适当的风险和交易限制】
    \begin{itemize}
    \item 也叫信贷额度（credit lines）
    \item 代表了公司愿意承担的最大风险金额（最大损失）
    \item 可对经常使用的交易对手、行业和国家附加限制。
    \item 因为在某些交易中，最大损失无法提前估计，而是取决于具体市场状况。所以，公司通常会使用不同限额的组合。
    \end{itemize}

\item D. Oversight：【配备具备独立性和资源的合格人员进行监督。】
    \begin{itemize}
    \item Independence 独立性
        \begin{itemize}
        \item 风险管理部门不应设在具有盈利中心的业务部门内。
        \item 风险经理的薪酬不应基于业务的盈利能力。
        \item 所有参与风险管理的员工都应该直接向首席风险官(CRO)汇报，是业务部门负责人。
        \item 首席风险官直接向首席执行官报告，并拥有接触董事会风险委员会和审计委员会的特权。
        \end{itemize}
    \item 资格要求 Qualifications：为了能更有效地执行其职责，风险管理需要具备专业知识和技能。
    \item Proximity to the Business Unit 与业务单位的接近：尽管风险管理部门需要保待独立性， 但与业务单位的接近能够使风险经理更好地了解业务， 并为业务提供支持和建议。
    \item Open Mind 开放的思维：一个好的风险经理应该具有开放的思维，能够提出建议以促成交易，而不仅仅是拒绝交易。
    \end{itemize}
\end{itemize}




\subsubsection{信贷委员会 Credit Committees}
    \begin{itemize}
        \item 最高的审批级别通常是信贷委员会（credit committee），由公司最高级别的管理人员组成。
        \begin{itemize}
            \item 信贷委员会由高级管理人员组成， 负责决策重要或高风险交易。
            \item  成员应包括组织的关键职能部门，如业务单位、风险管理、税务和会计、
            合规和法律。
            \item  委员会成员应具备足够的专业知识， 深刻理解风险。
            \item  委员会应有一名主席， 客观地收集意见并促进讨论。
            \item  委员会应有明确的章程， 详细说明审批流程。
            \item  一个运作良好的委员会应有批准和拒绝的交易。
            \item  会议记录应详细准确， 并在会议结束后能够迅速下发。
        \end{itemize}
    \end{itemize}
